\subsection{Validation using game variants}

To validate these agent sets, we elicit their response distributions to the costless and cycle versions of the game.
Both games maintain our fundamental requirement of similar but distinct data-generating processes. 
They involve reasoning well-explained by level-$k$ thinking, but feature payoff structures and incentives that differ from those of the basic game.
Most important, humans do not play these games in the same way as the basic version.
The forward KL divergence between the basic and costless game is $\KlBasicCostlessHuman$, and between the basic and cycle game is $\KlBasicCycleHuman.00$ (see Table~\ref{tab:AR_results} for the human response distributions for all three games).
These divergences are substantial: each is more than three times larger than the divergence between the human and optimized agent distributions in the right panel of Figure~\ref{fig:compare_agents_train}.
Accurate predictions on these distinct games would indicate that our optimized \aisub{s} capture generalizable patterns rather than merely replicating the original training scenario.

We elicit responses from all 1,000 \aisub{s} in our optimized sample $\bm{\theta}^*$ for both the costless and cycle versions of the game. As a benchmark for evaluating relative performance, we also elicit responses from the baseline AI 1,000 times per game. 
Figure~\ref{fig:compare_agents_test} presents these results, comparing the empirical distributions from AR's original human experiments (black lines) with those generated by the baseline AI (red) and the optimized mixture of theory-grounded \aisub{s} (blue).

\begin{figure}[h!]
\begin{center}
\begin{minipage}{\textwidth}
\caption{Response distributions for the cycle and costless versions of the 11-20 game} 
\label{fig:compare_agents_test}
\vspace*{-0.15in}
\centering
\includegraphics[width=.9\textwidth]{plots/cycle_costless_1120.pdf}
\end{minipage}
\end{center}
\begin{footnotesize}
\begin{singlespace}
\vspace*{-0.05in}
\emph{Notes:} This figure displays empirical PMFs for the costless (top row) and cycle (bottom row) variants of the 11-20 game. 
The columns correspond to the baseline (red), the optimized \aisub{s} (blue).
Within each panel, the empirical PMFs from \citeauthor{1120Arad2012} are imposed in black.
The KL divergence between each human and the AI response distribution is displayed in each panel.
For both variants of the game, the selected \aisub{s} (blue) are far closer to the human distribution than the baseline (red), even though the selected \aisub{s} are constructed using only the basic version of the game.
\end{singlespace}
\end{footnotesize}
\end{figure}

Consistent with the basic version of the game and \cite{gao2024caution}, the baseline AI overwhelmingly selects 19 shekels in both variants, a considerable divergence from AR's human subjects. 
In contrast, $\bm{\theta}^*$ is a far better predictor of both validation settings. 
The costless game in particular shows substantial improvement, with the KL divergence decreasing by \KlPercDiffCostless\% ($d(P, \hat P_{\bm\theta^*}) = \KlHumanOptCostlessDist$ vs. baseline $d(P,\hat P_0) = \KlHumanRawCostlessDist$).
It is almost perfectly predictive of the human responses when comparing the empirical PMFs.
In the cycle game, the KL divergence between the optimized AI and human responses is also reduced substantially, by \KlPercDiffCycle\% relative to the baseline ($d(P, \hat P_{\bm\theta^*}) = \KlHumanOptCycleDist$ vs. $d(P,\hat P_0) = \KlHumanRawCycleDist$).
Taken together, these results are interpreted as evidence that $\bm{\theta}^*$ has effectively generalized to the validation data---data that was not used to construct its mixture.
We therefore gain confidence that these agents may better predict entirely new settings that call for similar strategic reasoning.
